\section{The Tomb: Ostian Way, Basilica, Sarcophagus}

\subsection{Gaius's trophy and the burial tradition}

The burial tradition is as old, and as Roman, as the death tradition, and
it has never moved. Gaius, c.~200: the ``trophies of the apostles'' stand
``at the Vatican'' and ``on the Ostian way'' (via Eusebius, \work{Church
History} 2.25.7)---a monument one could walk out and see, in a suburban
cemetery zone along the road to Ostia, within about 130 years of the
death. Jerome, 392: ``buried in the Ostian way'' (\work{De viris
illustribus}~5). The later romances add the landowner: burial ``on the
estate of Lucina,'' a Christian matron (so the \work{Passion of Paul} of
Pseudo-Abdias, fourth century, as summarized in Benedict~XVI's catechesis
of 4 February 2009). The current official presentation of the basilica
makes the same claims in the same order, adding the archaeological
observation that the site lay in a genuine pagan-and-Christian necropolis,
and that a citizen's body could lawfully receive such a burial
(``The Tomb of the Apostle,'' basilica website via vatican.va, cited in
the References).

\subsection{From memorial to basilica}

Constantine enclosed the memorial in a first church---consecrated,
by the official account, on 18 November 324 by Pope Sylvester~I; because
the site was cramped between road and river, the emperors Valentinian~II,
Theodosius, and Arcadius replaced it from 386 with the immense five-aisled
``basilica of the three emperors'' (consecrated 395), which, uniquely
among Rome's great churches, preserved its ancient form until modern
times. On the night of 15--16 July
1823 fire destroyed it almost entirely; the rebuilt basilica, consecrated
by Pius~IX on 10 December 1854, reproduces the old plan, and the destruction-and-rebuilding is
the reason the shrine's continuity had to be re-demonstrated
archaeologically rather than simply seen. Under the papal altar a marble
slab of the fourth--fifth century, found bearing the inscription
\textsc{Paulo Apostolo Mart}\ (``to Paul, apostle and martyr''), lies
1.37~m below the present altar table, pierced by holes---one of them
connected by a small pipe to the tomb below, the ancient arrangement for
lowering cloths and pouring perfumes to touch the grave (basilica
presentation, as above).

\subsection{The sarcophagus examined, 2002--2009}

Beneath the slab stands a massive stone sarcophagus (2.55~m long by
1.25 by 0.97, per the official presentation), brought back to light in
excavations conducted by Vatican archaeologists between 2002 and 2006,
when a window was opened below the altar so that pilgrims can now see the
tomb. During the Pauline Year the Holy See authorized a minimally invasive
scientific examination, and Benedict~XVI announced its results in the
closing homily, at the tomb, on 28 June 2009:

\begin{studyframe}\small
``A tiny hole was drilled in the sarcophagus, which in so many centuries
had never been opened, in order to insert a special probe which revealed
traces of a precious purple-coloured linen fabric, with a design in gold
leaf, and a blue fabric with linen threads. Grains of red incense and
protein and chalk substances were also found. In addition, minute
fragments of bone were sent for carbon-14 testing by experts unaware of
their provenance. The fragments proved to belong to someone who had lived
between the first and second centuries. This would seem to confirm the
unanimous and undisputed tradition which claims that these are the mortal
remains of the Apostle Paul'' (Homily, Basilica of St.~Paul
Outside-the-Walls, 28 June 2009).
\end{studyframe}

\witnesslimit{The evidential ceiling must be stated with the
announcement. What the 2009 report establishes, at the reported level, is
that the venerated sarcophagus contains ancient burial furnishings and
human remains radiocarbon-consistent with a first--second-century death
---i.e., that the tomb is not a medieval confection. Radiocarbon cannot
identify a person; no full laboratory publication, calibration data, or
osteological study was located by this research; and the pope's own verb
is ``would seem to confirm.'' The identification of the remains as Paul's
rests, exactly as before 2009, on the continuity of the site tradition
from Gaius backward toward the burial---early, unrivaled, and unprovable.
This study affirms the tradition at that ceiling and no higher.}
